%% Preface
\chapter*{Preface}
\markboth{Preface}{Preface}
\addcontentsline{toc}{chapter}{Preface}

This book grew out of a correspondence. Over the course of 2026 Warren Ewens sent me a series of
questions, some about the Fundamental Theorem of Natural Selection, on which he had just written
what he called the end of a story, and more about the sampling formula that carries his name. What
happens to the formula when a Poisson Dirichlet subdivision is sampled with a nonuniform density?
How much of that survives in the two parameter family? What is the mean number of singleton colours
in a paintbox? Does a derivation from de Finetti's theorem say anything about probabilities in
Kingman's coalescent? Each question became a note, most of the notes became papers, and the papers
turned out to be chapters of a single argument.

The argument is that the classical questions of population genetics are statistical questions, and
that the statistical tools developed in recent decades (global local shrinkage priors,
Bayesian nonparametrics, data augmentation, and generative methods for inference without a
likelihood) answer some of them in ways that were not available to Fisher, Haldane, Wright, Kimura
or Ewens himself. In the other direction, several objects that statisticians now use routinely,
from the Dirichlet process to the horseshoe prior, turn out to be objects that population
geneticists defined first, for reasons of their own, or to share their marginal laws exactly.

The book has three parts, framed by a historical introduction, a chapter of mathematical
preliminaries, and a closing chapter on open problems. Chapter~\ref{ch:intro} tells the history of
the subject and introduces its basic models. Chapter~\ref{ch:toolkit}, \emph{A Mathematical
Toolkit}, collects in one place the probability and statistics that the later chapters draw on
repeatedly, so that those chapters can get to their arguments quickly. Part I
(Chapters~\ref{ch:ftns}, \ref{ch:hs} and \ref{ch:mp}) is about selection and Fisher's additive
model; Part II (Chapters~\ref{ch:esf}, \ref{ch:py}, \ref{ch:pig} and \ref{ch:rc}) about the Ewens
sampling formula and the sampling theory of neutral alleles; Part III (Chapters~\ref{ch:suff} and
\ref{ch:coal}) about inference for the coalescent when the likelihood cannot be computed.
Chapter~\ref{ch:future} collects the open problems, stated where possible as precise questions.
Each chapter has its own summary, notation and
references, so that it can be read independently. The cross references between chapters are there
for the reader who wants to follow a thread.

Two features are meant to make the book usable for teaching and for self study. Every chapter
contains worked examples, set apart from the text, in which the main results are carried out on
cases small enough to check by hand or with a few lines of code, and most chapters contain at least one
analysis of real genetic or ecological data. And every chapter ends with a set of exercises,
mixing proofs, calculations and short computations, each with a hint.

The chapters need not be read in order. Some suggested paths:
\begin{itemize}[leftmargin=*]
\item A course on the sampling formula: Chapters~\ref{ch:toolkit}, \ref{ch:esf}, \ref{ch:py}
and \ref{ch:pig}, adding Chapter~\ref{ch:rc} for a longer course.
\item A course on selection and the Fundamental Theorem: Chapters~\ref{ch:toolkit},
\ref{ch:ftns}, \ref{ch:hs} and \ref{ch:mp}.
\item A course on likelihood free inference: Chapters~\ref{ch:toolkit}, \ref{ch:esf},
\ref{ch:suff} and \ref{ch:coal}.
\item A short introduction for a reader new to the subject: Chapter~\ref{ch:intro}, then
Chapter~\ref{ch:toolkit}, then the first three sections of Chapter~\ref{ch:esf}, and finally
Chapter~\ref{ch:future}.
\end{itemize}
Chapter~\ref{ch:intro} can serve as background reading for any of these.

The prerequisites are a working knowledge of probability and of Bayesian inference at the level of
a first graduate course. No previous acquaintance with population genetics is assumed;
Chapter~\ref{ch:intro} supplies the history and the basic models, Chapter~\ref{ch:toolkit} the
common mathematical tools, and each later chapter defines what it uses. The numerical results and
worked examples in every chapter were produced by short Python programs, and the figures can be
reproduced from them.

\section*{Acknowledgements}

My first and largest debt is to Warren Ewens. This book exists because of our correspondence: a
long run of emails in which he asked the questions that started almost every chapter, read the
answers as they arrived, pointed out what was wrong or unclear in them, and then, more often than
not, asked the next question. I am grateful for every one of those emails, for the generosity with
which he shared his preprints and his memory of how the subject developed, and for the motivating
questions themselves, which were always better than the answers I first gave. His work over more
than half a century defines much of the territory the book covers, and Chapter~\ref{ch:intro}
tries to give some account of it. The book is dedicated to him.

I am also grateful to Vadim Sokolov for many discussions, in particular of double descent and of
overparameterised models, and to the Booth School of Business at the University of Chicago for
its support. The errors that remain are mine.

\vspace{1em}
\noindent Nicholas G.\ Polson\\
Chicago, October 2026
